% =====================================================================
%  THÈME 2 — Notation scientifique & unités — EXERCICES — Niveau MOYEN
% =====================================================================
\documentclass[11pt,a4paper]{article}
\usepackage[utf8]{inputenc}
\usepackage[T1]{fontenc}
\usepackage{lmodern}
\usepackage[french]{babel}
\usepackage{amsmath,amssymb,mathtools}
\usepackage{icomma}
\usepackage{siunitx}
\sisetup{output-decimal-marker={,},per-mode=symbol}
\usepackage[version=4]{mhchem}
\usepackage{xcolor}
\usepackage[most]{tcolorbox}
\usepackage{enumitem}
\usepackage{fancyhdr}
\usepackage[a4paper,margin=2.2cm,headheight=26pt,headsep=8pt]{geometry}

\definecolor{primary}{RGB}{150,98,162}
\definecolor{primarydark}{RGB}{92,52,110}
\newcommand{\cs}{chiffre significatif}
\newcommand{\CS}{chiffres significatifs}

\pagestyle{fancy}\fancyhf{}
\fancyhead[L]{\small\textcolor{primarydark}{\textbf{Fiche 1} — Chiffres significatifs}}
\fancyhead[R]{\small\textcolor{primarydark}{\textsc{Exercices · Moyen · Thème 2}}}
\fancyfoot[C]{\small\thepage}
\renewcommand{\headrulewidth}{0.8pt}
\renewcommand{\headrule}{\hbox to\headwidth{\color{primary}\leaders\hrule height \headrulewidth\hfill}}

\newtcolorbox[auto counter]{exo}[1][]{enhanced,breakable,boxrule=0.8pt,arc=2pt,
  colback=primary!4,colframe=primary,left=3mm,right=3mm,top=2mm,bottom=2mm,
  fonttitle=\bfseries,coltitle=white,
  title={Exercice~\thetcbcounter\ifx\relax#1\relax\relax\else\ ---\ \textmd{\itshape #1}\fi}}

\setlength{\parindent}{0pt}\setlength{\parskip}{3pt}

\begin{document}

\begin{center}
{\Large\bfseries\color{primarydark} Exercices — Niveau Moyen}\\[1pt]
{\color{primary} Thème 2 : notation scientifique \& changements d'unité}
\end{center}
\vspace{1mm}
\textit{On change d'unité en vérifiant à chaque fois que le nombre de \CS\ ne bouge pas, et on
apprend à repérer les nombres exacts.}
\vspace{2mm}

\begin{exo}[Une pesée dans plusieurs unités]
Une balance affiche $m = \qty{4,750}{\gram}$.
\begin{enumerate}[label=\alph*),leftmargin=1.6em,itemsep=1pt]
  \item Combien de \CS\ possède cette valeur ?
  \item Exprime $m$ en milligrammes, en kilogrammes, puis en tonnes.
  \item Vérifie que chacune de ces écritures possède le \emph{même} nombre de CS. Quel théorème
        du cours cela illustre-t-il ?
\end{enumerate}
\end{exo}

\begin{exo}[Un volume prélevé à la pipette]
On prélève $V = \qty{3,50}{\centi\litre}$ de solution.
\begin{enumerate}[label=\alph*),leftmargin=1.6em,itemsep=1pt]
  \item Exprime $V$ en millilitres, en litres et en microlitres.
  \item Donne l'écriture de $V$ en litres, en \emph{notation scientifique}.
  \item Combien de CS dans chaque écriture ? Conclus.
\end{enumerate}
\end{exo}

\begin{exo}[Étiquette d'eau minérale]
Une étiquette indique une teneur en ions calcium $[\ce{Ca^{2+}}] = \qty{90,0}{\milli\gram\per\litre}$.
\begin{enumerate}[label=\alph*),leftmargin=1.6em,itemsep=1pt]
  \item Exprime cette concentration en \si{\gram\per\litre}, puis en \si{\micro\gram\per\litre}.
  \item Écris la valeur en \si{\gram\per\litre} sous forme scientifique.
  \item Le nombre de CS ($3$) est-il modifié par ces conversions ? Justifie.
\end{enumerate}
\end{exo}

\begin{exo}[Exact ou mesuré ?]
Pour chaque nombre, dis s'il est \textbf{exact} (infinité de \CS) ou \textbf{mesuré} (nombre fini de CS).
\begin{enumerate}[label=\alph*),leftmargin=1.6em,itemsep=1pt]
  \item les $2$ atomes d'hydrogène de \ce{H2O} ;
  \item la masse $\qty{5,0}{\gram}$ lue sur une balance ;
  \item le facteur $\qty{1}{\hour}=\qty{3600}{\second}$ ;
  \item le nombre $\pi$ ;
  \item les $6$ faces d'un dé ;
  \item le volume $\qty{12,0}{\litre}$ mesuré au laboratoire.
\end{enumerate}
Question finale : un nombre exact peut-il \emph{limiter} le nombre de CS du résultat d'un calcul ?
\end{exo}

\begin{exo}[Convertir avec un facteur exact]
On convertit une masse $m = \qty{6,25}{\gram}$ ($3$~CS) en kilogrammes à l'aide de la relation
\emph{exacte} $\qty{1}{\kilo\gram}=\qty{1000}{\gram}$.
\begin{enumerate}[label=\alph*),leftmargin=1.6em,itemsep=1pt]
  \item Pose le calcul et donne le résultat en kilogrammes.
  \item Le $\num{1000}$ (exact) a-t-il fait perdre des CS au résultat ? Combien de CS conserve $m$ ?
\end{enumerate}
\end{exo}

\end{document}
